\begin{proof}

\begin{enumerate}[label={$\langle1\rangle$\arabic*.},leftmargin=*]
\item Let $T_0$ be the minimal Hopf surface. Every
irreducible curve $E$ on $T_0$ is smooth elliptic.
Indeed $b_2(T_0)=0$ gives $E^2=K_{T_0}\cdot E=0$,
so adjunction gives $p_a(E)=1$. A rational
normalization would define a nonconstant map
$\mathbb P^1\to T_0$. It lifts to the universal
cover $\mathbb C^2\setminus\{0\}$, where its two
coordinate functions are constant. This is
impossible. Thus the normalization genus is one
and the sum of the local delta invariants is zero.
Distinct curves are disjoint, since their
intersection number is zero and every local
intersection of distinct complex curves is
positive.

\item Push the effective conductor $D$ to $D_0$ on
$T_0$. The canonical comparison of P4.12 makes
$K_{T_0}+D_0$ torsion. The divisor $D_0$ is
nonempty: otherwise a positive power of $K_{T_0}$
would be trivial, contradicting the Kodaira
dimension $-\infty$ of a Hopf surface. Its
reduced support consists of $r\ge1$ disjoint
smooth elliptic curves. The blowup coefficient
$\mu-1$ calculation of P4R.3 applies without a
ruling. At a smooth conductor point it adds a
rational leaf precisely when $\mu\ge2$; at a
node it subdivides the corresponding edge.
Consequently $D_{\rm red}$ has $r$ connected
components, each a smooth elliptic root with a
rational tree attached. With $v$ rational
components it has exactly $v$ double points,
and hence $e(D_{\rm red})=v$. Since
$e(T)=k$, normalization away from the conductor
gives $e(S_{\rm reg})=k-v$. Each rational
component was created by one of the $k$ blowups,
so $k-v\ge0$.

\item Write $s$ for the number of connected components
of $C$, $c$ for its number of irreducible
components, and $g$ for the number of its
normalizations of genus one. Every normalized
component has genus zero or one: it receives a
finite nonconstant map from a normalized
conductor component. Each elliptic component
must receive an elliptic root, and distinct
image components require distinct roots.
Thus $c\ge g$ and $r\ge g$. Let $b$ be the sum
of the cycle ranks of the point-identification
graphs of the $s$ connected components. The
normalization formula gives
\[
 e(C)=c+s-2g-b.
\]

\item We retain the disconnected conductor in the
fundamental-group comparison. Its pushout has
one vertex for $T$ and one for each connected
component of $C$, with $r$ edges supplied by
the connected components of $D_{\rm red}$.
This graph is connected, because the finite
map $D_{\rm red}\to C$ is surjective.
The van Kampen presentation has $r-s$ stable
letters, in addition to the vertex groups.
The group of each conductor component is the
group of its elliptic root. Its map into $C$
factors, on that root, through the smooth
normalization of one image component; the
attached rational tree adds no fundamental
group generator. Kill the vertex group
$\pi_1(T)$ and the normal subgroups generated
by all normalized component groups of $C$.
Every conductor-group relation then disappears.
The resulting quotient is exactly
\[
 \pi_1(S)\twoheadrightarrow F_{b+r-s}.
\]
This is also the group of the homotopy-pushout
graph obtained by making the same quotients in
the normalization graph model of $C$. P4R.1
therefore gives $b+r-s\le1$. Substitution yields
\[
 e(C)\ge c+r-2g-1\ge-1.
\]

\item We exclude the only negative possibility by
the actual local hypersurface geometry.
The value $e(C)<0$ in these inequalities
forces $c=r=g$ and $b+r-s=1$. Every image
component then has elliptic normalization.
No rational conductor component can map
finitely onto one of them, by
Riemann--Hurwitz. Thus $D_{\rm red}$ consists
of disjoint smooth elliptic curves. This
conclusion permits blowups at coefficient-one
points; their exceptional curves are absent
from the conductor and do not affect the
following argument.

\item At a point $p\in C$, suppose there were two
distinct reduced curve germs $C_1,C_2$ of
$C$. Choose normalization preimages over $p$
whose conductor germs map onto $C_1,C_2$.
They must be distinct preimages, since the
reduced conductor has just one smooth germ
at each of its points. They correspond to
distinct irreducible analytic surface
branches $S_1,S_2$ of $S$ at $p$.
The support $S$ is a hypersurface in the
smooth total threefold. Its distinct
surface branches have a one-dimensional
intersection germ through $p$: two local
nonassociate irreducible equations in the
regular three-dimensional local ring have
height two. That intersection lies in the
conductor image of both branches. The
unique conductor germ on each chosen
normalization branch therefore has the
same image germ, giving $C_1=C_2$, a
contradiction.

\item Hence $C$ is locally unibranch and its
irreducible components are disjoint. Its
point-identification graphs have $b=0$
and $s=c=g$. The Euler formula gives
$e(C)=0$, contrary to the negative
possibility. We have proved $e(C)\ge0$,
and the asserted formulas finish the proof.

\end{enumerate}
\end{proof}
